\documentclass[11pt]{article}

\usepackage{amsmath, amssymb, amsthm}
\usepackage[margin=1in]{geometry}
\usepackage{hyperref}

\newtheorem{proposition}{Proposition}
\newtheorem{lemma}{Lemma}

\title{Rational Parametrization of the Hyperbola $X^2 - Y^2 = 1$}
\author{}
\date{}

\begin{document}

\maketitle

\section{Setup}

We seek all pairs $(X,Y) \in \mathbb{Q}^2$ satisfying
\begin{equation}
X^2 - Y^2 = 1.
\label{eq:main}
\end{equation}

The key observation is that the left-hand side factors as a difference of squares:
\begin{equation}
(X - Y)(X + Y) = 1.
\label{eq:factored}
\end{equation}

\section{Parametrization}

For any $t \in \mathbb{Q} \setminus \{0\}$, set
\begin{equation}
X + Y = t, \qquad X - Y = \frac{1}{t}.
\label{eq:sumdiff}
\end{equation}

This choice automatically satisfies \eqref{eq:factored}, since $(X-Y)(X+Y) = \frac{1}{t} \cdot t = 1$.

Solving the linear system \eqref{eq:sumdiff} for $X$ and $Y$ by adding and subtracting the two equations gives the parametrization:

\begin{equation}
\boxed{\,X(t) = \frac{1}{2}\left(t + \frac{1}{t}\right), \qquad Y(t) = \frac{1}{2}\left(t - \frac{1}{t}\right), \qquad t \in \mathbb{Q}\setminus\{0\}.\,}
\label{eq:param}
\end{equation}

\section{Completeness of the Parametrization}

\begin{proposition}
The map $t \mapsto (X(t), Y(t))$ defined in \eqref{eq:param} is a bijection between $\mathbb{Q}\setminus\{0\}$ and the set of rational points on the curve \eqref{eq:main}.
\end{proposition}

\begin{proof}
\emph{Every $t$ gives a solution.} Direct substitution into \eqref{eq:main} using \eqref{eq:param} confirms
\[
X(t)^2 - Y(t)^2 = (X(t)-Y(t))(X(t)+Y(t)) = \frac{1}{t}\cdot t = 1.
\]

\emph{Every solution arises this way.} Let $(X,Y) \in \mathbb{Q}^2$ satisfy \eqref{eq:main}. First note $X + Y \neq 0$: if $X+Y = 0$ then by \eqref{eq:factored} we would need $(X-Y)\cdot 0 = 1$, which is impossible. Hence we may set
\[
t := X + Y \in \mathbb{Q}\setminus\{0\}.
\]
From \eqref{eq:factored}, $X - Y = 1/t$. Solving this system for $X$ and $Y$ recovers exactly \eqref{eq:param}. Thus every rational solution corresponds to a unique nonzero rational $t$, namely $t = X+Y$.
\end{proof}

\section{Integer Form}

Writing $t = p/q$ in lowest terms, with $p, q$ nonzero integers and $\gcd(p,q) = 1$, the parametrization \eqref{eq:param} becomes

\begin{equation}
X = \frac{p^2 + q^2}{2pq}, \qquad Y = \frac{p^2 - q^2}{2pq}.
\end{equation}

\section{Geometric Interpretation}

This is the standard rational parametrization of a conic through a known rational point --- here $(X,Y) = (1,0)$. A line of slope determined by $t$ through this point meets the hyperbola in exactly one other rational point, and as $t$ ranges over $\mathbb{Q}\setminus\{0\}$ this sweeps out every rational point on the curve. The construction is directly analogous to the classical parametrization
\[
\left(\frac{1-t^2}{1+t^2}, \; \frac{2t}{1+t^2}\right)
\]
of the unit circle $X^2 + Y^2 = 1$; here the same idea is adapted to the hyperbola's factorization \eqref{eq:factored} rather than to the sum of squares.

\section{Verification}

Take $t = 2$:
\[
X = \frac{1}{2}\left(2 + \frac{1}{2}\right) = \frac{5}{4}, \qquad Y = \frac{1}{2}\left(2 - \frac{1}{2}\right) = \frac{3}{4}.
\]
Indeed,
\[
\left(\frac{5}{4}\right)^2 - \left(\frac{3}{4}\right)^2 = \frac{25}{16} - \frac{9}{16} = \frac{16}{16} = 1. \qquad \checkmark
\]

\end{document}
